% from ega3-6-fr.tex
\oldpage[III]{149}dérés;
nous le noterons donc
$\mathscr{T}\!or_n^S(\mathscr P_\bullet,\mathscr Q_\bullet)$; il résulte enfin
de (6.4.5) que pour toute partie ouverte $U$ (resp. $V$) de $X$ (resp. $Y$), la
restriction de $\mathscr{T}\!or_n^S(\mathscr P_\bullet,\mathscr Q_\bullet)$ à
$U\times_S V$ est canoniquement isomorphe à
$\mathscr{T}\!or_n^S(\mathscr P_\bullet|U,\mathscr Q_\bullet|V)$.
\end{env}

\begin{env}[6.5.3]
Passons maintenant au cas général où $S$ est quelconque, et soit $(S_\alpha)$
un recouvrement de $S$ formé d'ouverts affines; désignons par $X_\alpha$
(resp. $Y_\alpha$) l'image réciproque de $S_\alpha$ dans $X$ (resp. $Y$); il
faut encore prouver que les faisceaux
\[
  \mathscr{T}\!or_n^{S_\alpha}
  (\mathscr P_\bullet|X_\alpha,\mathscr Q_\bullet|Y_\alpha)
  =\mathscr G_\alpha
\]
vérifient la condition de recollement. Il suffit de définir, pour tout ouvert
affine $T$ contenu dans $S_\alpha\cap S_\beta$, des isomorphismes canoniques des
restrictions de $\mathscr G_\alpha$ et $\mathscr G_\beta$ à $U\times_S V$ (en
désignant par $U$ et $V$ les images réciproques de $T$ dans $X$ et $Y$
respectivement) sur
$\mathscr{T}\!or_n^T(\mathscr P_\bullet|U,\mathscr Q_\bullet|V)$; on peut, en
outre, se borner au cas où $T$ s'écrit à la fois $\mathrm D(f_\alpha)$ et
$\mathrm D(f_\beta)$, $f_\alpha$ (resp. $f_\beta$) étant une section de
$\mathscr O_S$ au-dessus de $S_\alpha$ (resp. $S_\beta$); mais alors
$\mathscr{T}\!or_n^T(\mathscr P_\bullet|U,\mathscr Q_\bullet|V)$ est
canoniquement isomorphe à
$\mathscr{T}\!or_n^{S_\alpha}(\mathscr P_\bullet|U,\mathscr Q_\bullet|V)$
d'une part, à
$\mathscr{T}\!or_n^{S_\beta}(\mathscr P_\bullet|U,\mathscr Q_\bullet|V)$
d'autre part, en vertu de (6.4.2); comme on vient de définir des isomorphismes
canoniques de $\mathscr G_\alpha$ sur
$\mathscr{T}\!or_n^{S_\alpha}(\mathscr P_\bullet|U,\mathscr Q_\bullet|V)$ et
de $\mathscr G_\beta$ sur
$\mathscr{T}\!or_n^{S_\beta}(\mathscr P_\bullet|U,\mathscr Q_\bullet|V)$
(6.5.2), cela achève de définir le $\mathscr O_Z$-Module
$\mathscr{T}\!or_n^S(\mathscr P_\bullet,\mathscr Q_\bullet)$. En outre, pour
toute partie ouverte $U$ (resp. $V$) de $X$ (resp. $Y$),
$\mathscr{T}\!or_n^S(\mathscr P_\bullet|U,\mathscr Q_\bullet|V)$ est
canoniquement isomorphe à la restriction de
$\mathscr{T}\!or_n^S(\mathscr P_\bullet,\mathscr Q_\bullet)$ à $U\times_S V$.

Il est immédiat que l'on a ainsi défini (dans les catégories de complexes de
Modules quasi-cohérents limités inférieurement) un bi-$\partial$-foncteur
$\mathscr{T}\!or_\bullet^S(\mathscr P_\bullet,\mathscr Q_\bullet)$ à valeurs
dans la catégorie des $\mathscr O_Z$-Modules, car il est clair que la question
est locale sur $X$, $Y$ et $S$, en vertu de (6.4.5) et de la remarque que
(6.4.4.2) est un morphisme de bi-$\partial$-foncteurs. On notera que si
$\mathscr P_\bullet$ et $\mathscr Q_\bullet$ sont réduits respectivement à
leurs termes de degré 0, $\mathscr F$ et $\mathscr G$,
$\mathscr{T}\!or_0^S(\mathscr F,\mathscr G)$ n'est autre, en vertu de
(6.4.1.1), que le produit tensoriel externe $\mathscr F\otimes_S\mathscr G$
défini dans (I, 9.1.2); cela résulte en effet de (I, 9.1.3).
\end{env}

\begin{env}[6.5.4]
Il résulte de la construction précédente et des remarques faites dans (6.4.6)
que $\mathscr{T}\!or_\bullet^S(\mathscr P_\bullet,\mathscr Q_\bullet)$ est
l'aboutissement de deux foncteurs spectraux,
$'\mathscr E(\mathscr P_\bullet,\mathscr Q_\bullet)$,
$''\mathscr E(\mathscr P_\bullet,\mathscr Q_\bullet)$, de termes $E_2$ égaux à
\[
  '\mathscr E_{pq}^2
  =\mathscr H_p\bigl(\mathscr{T}\!or_q^S
    (\mathscr P_\bullet,\mathscr Q_\bullet)\bigr)
  \tag{6.5.4.1}\label{III.6.5.4.1.fr}
\]
\[
  ''\mathscr E_{pq}^2
  =\bigoplus_{q'+q''=q}
   \mathscr{T}\!or_p^S
   \bigl(\mathscr H_{q'}(\mathscr P_\bullet),
         \mathscr H_{q''}(\mathscr Q_\bullet)\bigr).
  \tag{6.5.4.2}\label{III.6.5.4.2.fr}
\]
La suite spectrale (6.5.4.2) est toujours régulière; les deux suites spectrales
sont birégulières si $\mathscr P_\bullet$ et $\mathscr Q_\bullet$ sont limités
inférieurement. Un autre cas où les deux suites précédentes sont birégulières
est le suivant :
\end{env}

\begin{env}[6.5.5]
Nous dirons que sur un espace topologique $T$ un faisceau d'anneaux
$\mathscr A$ est \emph{de dimension cohomologique $\leqslant n$} si, pour tout
$t\in T$ l'anneau $\mathscr A_t$ est de dimension cohomologique $\leqslant n$;
on dira alors aussi que l'\emph{espace annelé} $(T,\mathscr A)$ est
\emph{de dimension cohomologique $\leqslant n$}. On dira qu'un faisceau
d'anneaux (resp. un espace annelé) est de dimension cohomologique \emph{finie}
s'il existe un entier $n$ tel qu'il soit de dimension cohomologique
$\leqslant n$. On notera que si les $\mathscr A_t$ sont des anneaux locaux
(commutatifs) noethériens,

